\section{Detector Pool, Reward Specification, and Non-degeneracy Framework}
\label{sec:methods}
This section specifies the frozen detector artifacts, the reward specification used by both learned controllers, and the pre-deployment pairwise non-degeneracy screening test. The last of these is stated as a theorem so that it can be applied to future learned selectors without re-derivation.

\subsection{Detector Pool}
Each dataset has four separately fitted configurations: depth-6 TinyDT; L2-regularized LightLR; 25-tree, depth-10 MedRF; and a 64/32/16 HeavyMLP. All use threshold 0.5 and training seed 11. Detector names do not assert a resource ranking; the physical measurements determine the empirical order. During TON-IoT physical experiments all models remain resident, and an action selects one model for the next batch without loading, fitting, or updating it. This design deliberately isolates the effect of controller decisions from model instantiation cost so that any observed switching cost is attributable to inference itself, not to hot-swapping detectors.

\subsection{Causal State Encoding}
In the original fixed-load setting, before window $t$ the controller constructs $s_t$ per~\eqref{eq:state}. The kernel temperature is smoothed as $\widetilde T_t=0.9\widetilde T_{t-1}+0.1T_t$, initialized at $T_0$. The threat observation $h_{t-1}$ is the arithmetic mean attack probability from the single detector selected in the preceding window; it is \emph{not} the preceding ground-truth attack fraction. Queue utilization is bounded to $[0,1]$, and the preceding arrival count is divided by 100 and bounded to $[0,2]$. Initial values are $h_0=0.5$, $\lambda_0/100=1$, and $a_0=\text{TinyDT}$.

Each of the first four state variables has three bins. Their two boundaries are $(45,70)$~$^{\circ}$C, $(1/3,0.7)$ for threat, $(1/3,2/3)$ for queue utilization, and $(0.9,1.1)$ for arrival ratio. Including four previous actions yields $3^4\times 4=324$ encoded states. Boundary equality enters the higher bin. The flat index is row-major in the listed order, with previous action changing fastest.

Corrected R2 instead encodes current arrivals $N_t$, raw FIFO carryover $Q_t^-$ before adding current arrivals, prior selected-path threat, and prior action. Arrival bins are $N_t\leq50$, $50<N_t\leq100$, and $N_t>100$; carryover bins are zero, 1--100, and above 100; threat bins are below $1/3$, $[1/3,0.7)$, and at least $0.7$. Prior action varies fastest, giving $3^3\times4=108$ states in arrival/carryover/threat/action order. Initial threat is 0.5 and prior action TinyDT. Temperature is safety-only; selected-path probabilities update threat after service.

The executed order within a window is: (i) state construction from causal telemetry and lagged context; (ii) action commitment and log; (iii) selected-detector inference; (iv) outcome log; and (v) lagged-state update for $t+1$. The information boundary of Section~\ref{sec:system} therefore holds at every decision; ground truth enters only offline training rewards, after the outcome is logged. This boundary is what allows preference-reversal analysis to be conducted purely on state features that the deployment can observe.

\subsection{Reward and Cost Provenance}
For the original fixed-load experiments, the unchanged per-window training reward is
\begin{equation}\label{eq:reward}
    r_t = 0.5\,B_t \;-\; 0.2\,\overline P_t \;-\; 0.2\,\overline L_t \;-\; 0.1\,\overline H_t,
\end{equation}
where $B_t$ is balanced accuracy over the classes present in the window. With $[x]_0^1=\min(1,\max(0,x))$, the normalization is
\begin{align}
    \overline P_t &= \biggl[\frac{P_t-2.347219336}{0.049806200}\biggr]_0^1,\\
    \overline L_t &= \biggl[\frac{L_t-1.2979595}{29.6547998}\biggr]_0^1,\\
    \overline H_t &= \biggl[\frac{\max(0,T_t^{\rm post}-\widetilde T_t)}{1\,^{\circ}\mathrm{C}}\biggr]_0^1.
\end{align}
$P_t$ and $L_t$ are frozen TON-IoT paced-profile lookups, not contemporaneous training measurements. Their means (power; batch latency) are 2.353132~W and 1.297960~ms for TinyDT, 2.347219~W and 1.544420~ms for LightLR, 2.397026~W and 30.952759~ms for MedRF, and 2.366688~W and 4.417836~ms for HeavyMLP. The compact software traces contain no post-action temperature response: their adapter supplies a fixed 35$^{\circ}$C predecision value and reuses that action-invariant value in the outcome field. Thus $\overline H_t=0$ throughout the executed training and validation traces; it is not an action-level thermal cost. No fitted thermal response or switching-energy penalty is used, and consequently no thermal control action is claimed.

Corrected R2 uses a separately frozen objective:
\begin{align}\label{eq:r2reward}
r_t^{\rm R2}={}&0.5 B_t-0.2\overline P_{a_t}-0.2\min(c_t,1)\nonumber\\
&-\ind[c_t>1\ \mathrm{or}\ Q_{t+1}>0]-\frac{Q_{t+1}}{N_t+Q_t^-},
\end{align}
with zero pending fraction for zero work. Training completion $c_t$ is $\lceil(N_t+Q_t^-)/100\rceil$ times the new warm-input p95 batch latency, in seconds. Power remains the historical TON min--max proxy, not contemporaneous variable-load energy. No thermal term is used. The fixed simulator develops no carryover, so queue ablation is not identifiable and is not executed; physical FIFO backlog is measured, not forced to zero.

\subsection{Pairwise Non-degeneracy Criterion}
Reward normalization is often reported as a hyperparameter without evidence that it permits any state to reverse the preferred action. We now show that this question can be answered before deployment with only the declared weights, cost profile, and a support region.

Consider two actions $a,b\in\mathcal M$. Decompose the scalar utility into a detection contribution and a normalized resource cost:
\begin{equation}
    U_a(s) = w_D\, D_a(s) - C_a,\qquad C_a=\sum_j w_j\,\overline C_{j,a},
\end{equation}
where $w_D=0.5$ in our specification and $\overline C_{j,a}$ collects the normalized resource penalties that are fixed on the declared support (as they are for the executed paced profiles at neutral thermal cost). Let $\Delta D_{ab}(s)=D_a(s)-D_b(s)$, $\Delta C_{ab}=C_a-C_b$, and $\tau_{ab}=\Delta C_{ab}/w_D$. Then
\begin{equation}\label{eq:diff}
    U_a(s)-U_b(s) = w_D\bigl(\Delta D_{ab}(s)-\tau_{ab}\bigr).
\end{equation}

\begin{theorem}[Pairwise non-degeneracy]\label{thm:nondeg}
Let $\mathcal S$ be a support region. Under the frozen scalar utility~\eqref{eq:diff}, a strict pairwise preference reversal between actions $a$ and $b$ exists on $\mathcal S$ if and only if
\begin{equation}\label{eq:nondeg}
    \min_{s\in\mathcal S}\Delta D_{ab}(s) < \tau_{ab} < \max_{s\in\mathcal S}\Delta D_{ab}(s).
\end{equation}
If \eqref{eq:nondeg} fails, the immediate difference cannot be strictly positive at one state and negative at another on $\mathcal S$. Thus an immediate-utility maximizer cannot reverse the pair there. This does not constrain a long-horizon RL policy when continuation values differ.
\end{theorem}

\begin{corollary}[Attainability under bounded detection]\label{cor:bound}
If $D_a(s),D_b(s)\in[0,1]$ for every $s\in\mathcal S$, then reversal is impossible whenever $|\tau_{ab}|\ge 1$; and reversal is guaranteed impossible whenever the attainable interval $[\min\Delta D_{ab},\max\Delta D_{ab}]$ does not straddle $\tau_{ab}$.
\end{corollary}

Because balanced accuracy lies in $[0,1]$, the corollary gives an immediate ceiling. Under the executed candidate-range scaling, the frozen nonthermal penalties imply $\tau_{\text{MedRF},\text{LightLR}}=0.7967$. A reproducible audit of all 2{,}500 windows in the five hash-bound validation traces finds $\Delta D_{\text{MedRF},\text{LightLR}}\in[-0.1053,0.1053]$, so this threshold lies outside the observed validation-window support. This retrospective audit excludes MedRF/LightLR and HeavyMLP/LightLR reversal on observed validation support. TinyDT/LightLR can reverse: TinyDT has higher immediate reward in 7 of 2{,}500 validation windows. These label-informed window comparisons establish neither optimality over all encoded states nor optimal long-horizon return, and do not fully explain the learned collapse. There is no evidence that this audit preceded test inspection. Supplementary Theorem~S1 gives the proof, Supplementary Table~S11A gives the full pairwise matrix, and Supplementary Algorithm~S1 gives the validation-only reward audit. If reversal is unattainable, the engineer should use an audited static schedule or renormalize against an external deployment contract---for example idle-relative power, the latency SLO, an energy envelope, or thermal headroom---rather than against detector-pool extrema.

The value $\tau\simeq0.7967$ belongs to the original fixed-cost reward only. Corrected R2 has load-dependent completion and service penalties, so this number is not transferred to its objective. Neither the theorem nor the corollary proves inevitable static RL behavior.

\subsection{Policies}
Four policy families are evaluated in the primary physical campaign. \emph{Static-\{TinyDT, LightLR, MedRF, HeavyMLP\}} freeze a single detector. \emph{CFSM} is a fixed heuristic initialized to TinyDT; at a smoothed temperature of at least 70~$^{\circ}$C it steps down one model rank, bypassing cooldown; otherwise it holds for two windows after a switch and steps down/up one rank when threat crosses 0.7. Queue and load are logged but not used by CFSM. \emph{Tabular-Q} uses zero initialization, $\alpha=0.1$, $\gamma=0.9$, and the update
\begin{equation}
    \begin{aligned}
        Q(s_t,a_t) &\leftarrow Q(s_t,a_t) + 0.1\,\delta_t,\\
        \delta_t &= r_t + 0.9\max_a Q(s_{t+1},a) - Q(s_t,a_t),
    \end{aligned}
\end{equation}
with zero terminal bootstrap~\cite{Watkins1992}. \emph{DQN} uses one-hot state, two 64-unit ReLU layers, Adam $10^{-3}$, replay 10{,}000, batch 32, and a target copy every 250 updates~\cite{Mnih2015}. In the original fixed-load setting both learners train for 1{,}000 episodes with $\varepsilon$ linearly annealed from 1 to 0.01; checkpoints are selected every 25 episodes on five validation traces. Test execution freezes the policy and performs no update.

Corrected R2 trains 25 artifacts using seeds 2003, 2011, 2017, 2027 and 2029: Q full/minus-arrival/minus-threat/minus-previous-action, and full-state DQN. Original learning hyperparameters are retained, but every run fixes its last checkpoint after 1,000 episodes of 500 windows. Full-state seed selection maximizes summed validation reward, ties choosing the lowest seed: Q 2011 and DQN 2003. The ten original training and five validation trace seeds are retained; validation resamples its existing pool. Historical test results were already viewed before this follow-up design. No physical result changes selection. Original checkpoint curves and training-time estimates do not describe these new runs.

Two causal policies are added in the R1/R2 campaigns (Section~\ref{sec:eval}). \emph{ConfidenceCascade} executes LightLR first and invokes MedRF on rows whose posterior lies in $[0.3,0.7]$; R1 uses 100-row windows, while corrected R2 serves FIFO batches of at most 100 rows. This post-entry confidence is not available to RL before window action commitment. In corrected R2, \emph{DeadlineGuard} selects MedRF when $\lceil(N_t+Q_t^-)/100\rceil L_{95}^{\rm RF}\leq1$~s and LightLR otherwise, using the same current arrivals and pre-append carryover exposed to RL. Its fixed calibration is not retuned using physical outcomes.

Archived process-start and policy-creation timestamps indicate approximately 26.7~s for the Tabular-Q training run and 256.9~s for DQN; these are provenance-derived wall-clock intervals rather than calibrated compute benchmarks, and training energy was not instrumented. Deployment cost is measured separately in Section~\ref{sec:dynamicresults}.
